% ASTRO living research paper.  arXiv-friendly: standard packages, no shell escape.
\documentclass[10pt,twocolumn]{article}

\usepackage[margin=0.72in]{geometry}
\usepackage{amsmath,amssymb,mathtools}
\usepackage{booktabs}
\usepackage{graphicx}
\usepackage{xcolor}
\usepackage{microtype}
\usepackage{enumitem}
\usepackage[numbers,sort&compress]{natbib}
\usepackage[hidelinks]{hyperref}
\usepackage{url}
\usepackage{caption}

\definecolor{astropurple}{HTML}{7A3E9D}
\definecolor{lightpurple}{HTML}{F5EFF8}
\definecolor{muted}{HTML}{666666}
\setlength{\columnsep}{0.24in}
\setlength{\parindent}{1em}
\setlength{\parskip}{0pt}
\captionsetup{font=small,labelfont=bf,skip=4pt}
\setlist{nosep,leftmargin=*}
\emergencystretch=1em

\newcommand{\ASTRO}{\textsc{Astro}}
\newcommand{\ASTROv}{\textsc{Astro-v2}}
\newcommand{\Muon}{\textsc{Muon}}
\newcommand{\NorMuon}{\textsc{NorMuon}}
\newcommand{\AdaMuon}{\textsc{AdaMuon}}
\newcommand{\norm}[1]{\lVert #1 \rVert}
\newcommand{\obs}[1]{%
  \medskip\noindent\fcolorbox{astropurple!40}{lightpurple}{%
  \parbox{\dimexpr\columnwidth-2\fboxsep-2\fboxrule}{%
  \textbf{Observation.} #1}}\medskip}
\newcommand{\pendingfig}[2]{%
  \fbox{\parbox[c][1.55in][c]{0.95\linewidth}{\centering\small\color{muted}%
  \textbf{#1}\\[3pt]#2}}}
\newcommand{\maybegraphic}[3][\linewidth]{%
  \IfFileExists{../../artifacts/figures/#2.pdf}%
    {\includegraphics[width=#1]{../../artifacts/figures/#2.pdf}}%
    {\pendingfig{Figure pending}{#3}}}

\title{\vspace{-1.2em}\textbf{ASTRO: Auditing and Improving Matrix-Structured\\Optimizers for Language Model Training}}
\author{Pop Alexandru\\\small Independent Researcher\\\small \url{https://github.com/pop123-ux/ASTRO}}
\date{September 2026 \;---\; living research draft}

\begin{document}
\maketitle
\vspace{-1.1em}

\begin{abstract}
Matrix-structured optimizers such as \Muon, \NorMuon, and \AdaMuon improve
language-model training by transforming matrix-valued updates rather than
treating every parameter independently. We document the development of \ASTRO,
a Muon-family optimizer and an experimental framework for asking which apparent
optimizer improvements survive fair tuning, longer horizons, larger models,
independent seeds, and wall-clock accounting.

The current \ASTROv recipe keeps Muon-style spectral updates, splits fused Q/K/V
blocks, applies a norm-preserving neuron-wise second moment after the polar step,
and routes non-operator parameters through an Adam-like scalar path. It also
removes two choices that ablation showed were harmful or accidental: the scalar
path no longer inherits spectral-path $\beta_1=0.95$, and cautious weight decay
is disabled. At GPT-2 124M and 900 steps, a targeted three-configuration study
gives mean paired validation-loss differences versus Muon of \textbf{-0.2208}
for \ASTROv, \textbf{-0.2137} for the beta-aligned ASTRO-MB variant, and
\textbf{-0.2047} when post-polar variance power is removed. All three variants
win all three shared configurations, but configurations are not independent
seeds. \ASTROv costs about \textbf{6.4\%} more wall-clock per 900-step run.

The evidence is promising but mixed. In an independently tuned 300-step 124M
benchmark, the earlier/plain ASTRO recipe is \textbf{+0.0180} validation loss
worse than Muon. The current horizon study is complete only at 300 and 600
steps, while independent-seed replication and the first 355M test are in
progress. A second line of evidence shows that fused spectral treatment can
allocate update mass strongly toward V, whereas splitting Q/K/V substantially
balances it. We therefore frame the current hypothesis narrowly: \ASTROv may
primarily improve \emph{robustness to hyperparameter choice}, rather than
uniformly shifting the best achievable loss. The paper preserves negative runs
and withdrawn explanations as first-class evidence rather than hiding the path
to the current recipe.
\end{abstract}

\section{Introduction}

Muon made a simple idea operational: hidden weight matrices should be optimized
as matrices, not merely as bags of coordinates \citep{jordan2024muon}. Later
methods add adaptivity after orthogonalization, notably \NorMuon
\citep{li2025normuon} and \AdaMuon \citep{si2025adamuon}. At the same time,
large comparative studies show that optimizer rankings depend strongly on
hyperparameter tuning, model scale, data budget, and the training horizon at
which performance is read out \citep{wen2025fantastic}.

These facts make new optimizer results easy to overstate. A method can appear
better because its search range is better centered, because an auxiliary
Adam-like path was tuned differently, because training ended at a favorable
intermediate point, or because the method simply performs more computation per
step. ASTRO began as an attempt to add additional control to Muon-family
updates. The project changed direction as experiments invalidated parts of the
original explanation.

This paper therefore makes two contributions. First, we study the current
\ASTROv recipe against modern Muon-family baselines. Second, we treat the
\emph{research process itself} as part of the artifact: shared configurations
are distinguished from shared seeds, wall-clock cost accompanies step-matched
loss, negative experiments remain visible, and incomplete cells stay visibly
incomplete in both plots and prose.

\paragraph{Contributions.}
\begin{enumerate}
  \item A PyTorch implementation of ASTRO that subclasses
  \texttt{torch.optim.Optimizer} for state management while implementing the
  matrix update explicitly.
  \item A 124M comparison to Muon, NorMuon, and AdaMuon that documents both the
  optimizer journey and the current ASTRO-v2 recipe.
  \item An ablation showing that ASTRO-MB captures almost all of the current
  ASTRO-v2 gain, while setting the post-polar variance power to zero preserves
  most of it.
  \item A QKV allocation analysis showing that fused spectral treatment can
  produce strongly imbalanced update allocation, while block-wise treatment
  restores substantially more balanced Q/K/V allocation.
  \item A reproducible paper pipeline: committed JSON is the evidence layer,
  Matplotlib generates publication figures, W\&B is optional, and the LaTeX
  source renders visible placeholders for missing measurements.
\end{enumerate}

\section{Background}

\subsection{Muon}

For a matrix gradient $G_t$, Muon maintains momentum and applies a short
Newton--Schulz-style iteration that approximates the polar factor of the
momentum matrix. If $M=U\Sigma V^\top$, the idealized transform is $UV^\top$.
Practical Muon uses this matrix path for hidden operators and an AdamW-like
auxiliary path for embeddings, biases, norms, and other non-matrix parameters.

The numerical subroutine is itself an active research topic. Polar Express
shows that changing the polynomial across iterations can improve approximation
quality while remaining GPU-friendly \citep{amsel2025polar}.

\subsection{NorMuon and AdaMuon}

NorMuon adds one adaptive second-moment statistic per output neuron after
orthogonalization and preserves the total update norm
\citep{li2025normuon}. AdaMuon applies element-wise adaptivity to the
orthogonalized update and aligns the RMS of the resulting direction
\citep{si2025adamuon}. They are important baselines because they test whether
post-spectral adaptation improves only peak loss or also the region of stable
hyperparameters.

\subsection{Why horizon and scale matter}

Wen et al. compare optimizers across model sizes and data/model ratios and show
that fair conclusions require optimizer-specific tuning and evaluation at the
intended final budget; rankings can change at intermediate checkpoints and
reported speedups shrink with scale \citep{wen2025fantastic}. We therefore treat
124M/300, 124M/900, and 355M as separate empirical questions.

Recent curvature work provides a complementary view: at matched validation
loss, Muon's advantage over Adam is associated with lower normalized
directional sharpness rather than simply smaller update norms
\citep{wang2026curvature}. Our current paper does not yet claim a curvature
explanation for ASTRO-v2; the repository contains instrumentation for such a
follow-up.

\section{Current ASTRO}

\subsection{Routing}

The current reference implementation is \texttt{class Astro(torch.optim.Optimizer)}
in \texttt{scripts/astro\_lab.py}. Subclassing PyTorch's optimizer base supplies
parameter groups, state, serialization, and interoperability; it does not supply
the ASTRO update.

Parameters are routed structurally. Genuine hidden 2-D operators take the
spectral path. Embeddings, the tied output head, biases, normalization
parameters, and other non-operators take an Adam-like scalar path. HuggingFace
GPT-2 \texttt{Conv1D} weights are transposed into operator orientation. Fused
attention projections can be split into Q, K, and V blocks before the spectral
transform.

\subsection{Spectral path}

ASTRO maintains momentum
\begin{equation}
M_t = \beta_1 M_{t-1} + (1-\beta_1)G_t,
\end{equation}
forms a Nesterov-style lookahead, and applies the approximate polar transform
block-by-block to obtain $Q_t$. It then tracks a row-wise second moment on the
\emph{orthogonalized} direction,
\begin{equation}
v_t = \beta_2 v_{t-1} + (1-\beta_2)\operatorname{rowmean}(Q_t^2),
\end{equation}
and redistributes the direction according to
\begin{equation}
\widetilde Q_t = Q_t \oslash \widehat v_t^{\gamma/2}.
\end{equation}
The original Frobenius norm is restored,
\begin{equation}
D_t = \widetilde Q_t\frac{\norm{Q_t}_F}{\norm{\widetilde Q_t}_F},
\end{equation}
so this stage mainly redistributes update energy rather than changing the total
step magnitude.

\subsection{ASTRO-v2 is a recipe}

ASTRO-v2 is not a second optimizer class. The harness instantiates the same
\texttt{Astro} implementation with
\begin{verbatim}
"astro_v2": {
  "betas": (0.9, 0.95),
  "cautious_wd": False,
}
\end{verbatim}
The earlier recipe used $\beta_1=0.95$ on both spectral and scalar groups, even
though Muon's Adam-like auxiliary path uses $0.9$. This meant a substantial
fraction of GPT-2 could differ before any matrix geometry was considered.
ASTRO-v2 removes this asymmetry and disables cautious weight decay.

ASTRO-MB isolates the beta alignment while retaining cautious weight decay.
ASTRO-v2-$\gamma0$ additionally disables the post-polar variance-power effect.
These variants are the main attribution controls in the current study.

\section{Experimental protocol}

\subsection{Model, data, and validation}

The scale campaign trains GPT-2-style causal language models from scratch on
FineWeb-Edu with GPT-2 BPE. The principal model has 12 layers, 12 heads, width
768, and approximately 124M parameters. Sequence length is 512. The first scale
check uses a 355M configuration.

Validation is pinned to the same 400,000-token slice at every training horizon.
This corrects an earlier protocol in which changing the maximum training budget
could also change the validation text.

\subsection{Two different kinds of pairing}

We use two distinct experimental units:
\begin{itemize}
  \item \textbf{shared configurations}: both optimizers receive the same
  learning rate, weight decay, and scalar-path multiplier; this measures
  sensitivity and robustness;
  \item \textbf{shared seeds}: frozen configurations are rerun with identical
  stochastic seeds; this measures stochastic reproducibility.
\end{itemize}
Three configurations are not three seeds, and the paper never uses one as a
surrogate for the other.

\subsection{Compute accounting}

A step-matched win is not automatically a compute win. We therefore report
wall-clock per run and prepare frozen validation trajectories for both
loss-vs-step and loss-vs-time comparisons. In the current targeted 124M/900
experiment, ASTRO-v2 costs about 6.4\% more wall-clock than Muon.

\section{Current results}

\subsection{Negative evidence at 300 steps}

The expanded tuned 124M/300 benchmark gives Table~\ref{tab:300}.

\begin{table}[t]
\centering
\caption{GPT-2 124M, 300 steps, three evaluation seeds. Lower is better. The
ASTRO row is the earlier/plain recipe, not ASTRO-v2.}
\label{tab:300}
\small
\begin{tabular}{lrr}
\toprule
Optimizer & Mean val. loss & $\Delta$ vs Muon \\
\midrule
NorMuon & \textbf{6.6280} & \textbf{-0.0077} \\
Muon & 6.6357 & 0.0000 \\
Earlier ASTRO & 6.6538 & +0.0180 \\
AdamW & 6.8445 & +0.2087 \\
\bottomrule
\end{tabular}
\end{table}

\obs{The ASTRO-v2 story cannot be back-projected onto earlier ASTRO runs: at
124M/300, the earlier recipe is worse than Muon under the expanded tuned
benchmark.}

\subsection{The optimizer journey at 124M/900}

At one shared configuration, the modern Muon-family comparison is shown in
Table~\ref{tab:shared} and Figure~\ref{fig:journey}.

\begin{table}[t]
\centering
\caption{One shared GPT-2 124M/900 configuration:
$\eta=0.010206$, weight decay $0.001060$, scalar multiplier $0.4369$. This is a
pointwise comparison, not a tuned global ranking.}
\label{tab:shared}
\small
\begin{tabular}{lrr}
\toprule
Optimizer & Val. loss & $\Delta$ vs Muon \\
\midrule
\textbf{ASTRO-v2} & \textbf{5.9208} & \textbf{-0.0827} \\
ASTRO-MB & 5.9226 & -0.0809 \\
AdaMuon & 6.0028 & -0.0007 \\
Muon & 6.0035 & 0.0000 \\
NorMuon & 6.0155 & +0.0121 \\
\bottomrule
\end{tabular}
\end{table}

\begin{figure*}[t]
\centering
\maybegraphic[0.96\textwidth]{fig_optimizer_journey}{Run
\texttt{python scripts/figures/make\_all.py --only optimizer\_journey}.}
\caption{\textbf{The optimizer journey.} Left: Muon, contemporary adaptive Muon
variants, and the current ASTRO variants at one shared configuration. Right:
paired deltas versus Muon across the three targeted mechanism settings. The
aggressive-LR configuration creates the largest separation, motivating a
robustness hypothesis rather than a universal-ranking claim.}
\label{fig:journey}
\end{figure*}

\obs{At this shared 124M/900 setting, almost the entire separation from Muon is
already present in ASTRO-MB; ASTRO-v2 improves on ASTRO-MB by only about 0.002
validation loss.}

\subsection{Targeted mechanism study}

Across three shared settings, the mean paired deltas versus Muon are shown in
Table~\ref{tab:mechanism}.

\begin{table}[t]
\centering
\caption{Targeted 124M/900 mechanism experiment. Settings are shared
hyperparameter configurations, not independent seeds.}
\label{tab:mechanism}
\small
\begin{tabular}{lrr}
\toprule
Variant & Mean $\Delta$ vs Muon & Settings won \\
\midrule
\textbf{ASTRO-v2} & \textbf{-0.2208} & \textbf{3/3} \\
ASTRO-MB & -0.2137 & 3/3 \\
ASTRO-v2 ($\gamma=0$) & -0.2047 & 3/3 \\
\bottomrule
\end{tabular}
\end{table}

The largest gap occurs at $\eta\approx0.0505$, where Muon reaches 6.4527 while
ASTRO-v2 remains at 5.9436. At the milder settings, the gaps are substantially
smaller.

\obs{Current evidence points more strongly toward \textbf{hyperparameter
robustness} than toward the post-polar variance-power term as the main source of
ASTRO-v2's advantage.}

\subsection{Compute cost}

\begin{figure}[t]
\centering
\maybegraphic{fig_efficiency_v2}{Run
\texttt{python scripts/figures/make\_all.py --only efficiency\_v2}.}
\caption{Validation loss against measured wall-clock for the three targeted
settings. Thin connectors join optimizers evaluated at the same hyperparameter
configuration. Lower-left is better.}
\label{fig:efficiency}
\end{figure}

Muon takes about 787 seconds per 900-step run in the targeted study; ASTRO-v2
about 838 seconds, a mean overhead of roughly 6.4\%.

\obs{ASTRO-v2's current result is a better-loss-at-equal-steps result. A full
time-matched or time-to-target experiment is required before calling it a
training-efficiency improvement.}

\subsection{Horizon: incomplete by design}

\begin{table}[t]
\centering
\caption{Current ASTRO-v2 horizon state. The 900 and 2700 cells are incomplete
and must not be interpreted as final averages.}
\label{tab:horizon}
\small
\begin{tabular}{lrrr}
\toprule
Steps & Completed & Mean $\Delta$ & Wins \\
\midrule
300 & 5/5 & -0.0500 & 2/5 \\
600 & 5/5 & -0.0541 & 3/5 \\
900 & \textbf{2/5} & -0.0118 so far & 1/2 \\
2700 & \textbf{0/5} & pending & pending \\
\bottomrule
\end{tabular}
\end{table}

\begin{figure}[t]
\centering
\maybegraphic{fig_horizon_v2}{Finish the horizon state or render the current
partial figure with \texttt{make\_all.py --only horizon\_v2}.}
\caption{Paired ASTRO-v2 minus Muon validation-loss deltas through training.
Filled diamonds denote complete cells; an open diamond denotes a partial cell.
Individual points are shared configurations.}
\label{fig:horizon}
\end{figure}

One favorable configuration stays near $-0.09$ from 300 through 900 steps. A
second configuration remains unfavorable, so the effect is not uniform across
the search space.

\obs{The current horizon data are consistent with a persistent effect at some
configurations, but they do not yet determine the final 900- or 2700-step
average trend.}

\subsection{Training trajectories}

\begin{figure*}[t]
\centering
\maybegraphic[0.96\textwidth]{fig_training_trajectories}{Generate frozen
validation trajectories with \texttt{scripts/paper/run\_trajectory.py}.}
\caption{Planned final learning-dynamics figure. Left: validation loss versus
training step. Right: the same trajectories versus measured training wall-clock
with validation-probe time removed. These runs are generated only after
configurations are frozen.}
\label{fig:trajectories}
\end{figure*}

This figure is intentionally allowed to remain a visible placeholder in a living
draft. The paper build must never reuse an old trajectory merely because a new
one is absent.

\section{Fused QKV allocation}

A fused attention projection stores Q, K, and V in one matrix. A spectral
transform can be mathematically valid for that fused matrix while distributing
update energy unevenly among the three operators inside it.

\begin{table}[t]
\centering
\caption{Example Q/K/V update allocation. Splitting is applied before the
spectral transform.}
\label{tab:qkv}
\small
\begin{tabular}{lcc}
\toprule
Checkpoint & Fused Q/K/V & Split Q/K/V \\
\midrule
GPT-2 & .244/.226/\textbf{.530} & .318/.347/.336 \\
GPT-2 Medium & .245/.266/\textbf{.489} & .308/.369/.323 \\
Pythia-410M & .218/.248/\textbf{.534} & .356/.300/.344 \\
\bottomrule
\end{tabular}
\end{table}

The corresponding gradient V/Q ratios are approximately $2.65\times$,
$3.27\times$, and $3.28\times$. Early GPT-2 layers can be more extreme. Random
initialization also exhibits V-skew, so training is not claimed as the sole
cause.

\obs{Fused spectral treatment can produce strongly imbalanced Q/K/V allocation,
while block-wise treatment restores substantially more balanced allocation. The
imbalance is architecture- and stage-dependent rather than purely learned during
training.}

Splitting itself has a cost: earlier T4 timing in this repository shows that
block-wise orthogonalization becomes materially more expensive at larger widths.
This motivates testing a scheduled split instead of assuming that a permanent
split is always optimal.

\section{Research journey and corrections}

\subsection{A shared learning rate was not a fair comparison}

Early runs forced Muon and ASTRO to share one learning rate even though their
spectral and scalar update scales were not calibrated identically. Opposite
conclusions appeared at different fixed rates. This motivated optimizer-specific
tuning and the present shared-configuration design.

\subsection{The scalar path mattered more than expected}

ASTRO originally used $\beta_1=0.95$ for both matrix and scalar groups. Muon's
auxiliary Adam path uses $0.9$. Aligning that scalar behavior produced a larger
improvement than several more elaborate spectral changes.

\subsection{Cautious weight decay was not a free safety mechanism}

A masked weight decay deletes decay where the mask fails. Unless survivors are
rescaled, the method silently changes effective regularization. ASTRO-v2 disables
this mechanism rather than treating it as a harmless default.

\subsection{Variance power is not yet the headline mechanism}

Setting $\gamma=0$ loses only about 0.016 validation loss on average relative to
full ASTRO-v2 in the targeted run while retaining a large advantage versus Muon.
The most mathematically distinctive component is therefore not automatically the
component carrying the result.

The repository retains \texttt{artifacts/measured.json} as a historical ledger
and uses the smaller \texttt{artifacts/paper\_results.json} as the current
paper-facing snapshot.

\section{What must be completed before submission}

\paragraph{Horizon.} Complete the remaining 900-step settings and all five
2700-step settings.

\paragraph{Independent seeds.} Freeze configurations first, then evaluate shared
seeds. In the current CLI, \texttt{--seeds 3} means seeds 100, 101, and 102;
\texttt{--seeds 100} means one hundred seeds.

\paragraph{Scale.} The first 355M/900 shared-configuration comparison is the
cheapest falsification of a small-model-only effect.

\paragraph{Time matching.} Generate frozen validation trajectories and report
both equal-step and equal-wall-clock results, preferably including time-to-target.

\section{Limitations}

The strongest current ASTRO-v2 result is a configuration-paired result, not yet
an independent-seed result. The principal scale is only 124M. The current v2
horizon is unfinished. ASTRO-v2 is measurably more expensive per step than Muon.
Some historical experiments used superseded validation or tuning protocols and
remain in the repository only as an audit trail. Finally, although ASTRO-MB's
near-tie with ASTRO-v2 narrows the mechanism hypothesis, more isolation is needed
before claiming that beta alignment alone is causal.

\section{Conclusion}

ASTRO-v2 is currently \textbf{promising, not proven}. At GPT-2 124M/900 it
substantially improves on Muon across three shared configurations, with the
largest separation in an aggressive-learning-rate regime. Yet the evidence also
says what the result is not: the earlier ASTRO recipe loses in a tuned 300-step
benchmark, post-polar variance power does not explain most of v2's current gain,
and the step-matched improvement costs roughly 6\% additional wall-clock.

The emerging hypothesis is therefore narrower and more testable: ASTRO-v2 may
improve the \emph{robustness} of matrix-structured optimization to hyperparameter
choice, while a simpler subset of its recipe may carry most of the effect. The
ongoing horizon, seed, scale, and time-matched experiments are designed to
falsify that hypothesis quickly.

The methodological commitment is simpler: the final paper will show the journey,
not just the winner. Negative results, protocol corrections, and unfinished
cells remain visible because optimizer research is most useful when the
comparison can fail.

\paragraph{Reproducibility.}
The complete manuscript, experiment ledger, current paper snapshot, plotting
code, and build script are available at
\url{https://github.com/pop123-ux/ASTRO}.

\begin{thebibliography}{99}
\bibitem[Jordan et al.(2024)]{jordan2024muon}
K.~Jordan, Y.~Jin, V.~Boza, J.~You, F.~Cesista, L.~Newhouse, and J.~Bernstein.
\newblock Muon: An optimizer for hidden layers in neural networks.
\newblock 2024. \url{https://kellerjordan.github.io/posts/muon/}.

\bibitem[Li et al.(2025)]{li2025normuon}
Z.~Li et al.
\newblock NorMuon: Making Muon more efficient and scalable.
\newblock \emph{arXiv:2510.05491}, 2025.

\bibitem[Si et al.(2025)]{si2025adamuon}
C.~Si, D.~Zhang, and W.~Shen.
\newblock AdaMuon: Adaptive Muon Optimizer.
\newblock \emph{arXiv:2507.11005}, 2025.

\bibitem[Wen et al.(2025)]{wen2025fantastic}
K.~Wen, D.~Hall, T.~Ma, and P.~Liang.
\newblock Fantastic Pretraining Optimizers and Where to Find Them.
\newblock \emph{arXiv:2509.02046}, 2025.

\bibitem[Amsel et al.(2025)]{amsel2025polar}
N.~Amsel, D.~Persson, C.~Musco, and R.~Gower.
\newblock The Polar Express: Optimal Matrix Sign Methods and Their Application to
  the Muon Algorithm.
\newblock \emph{arXiv:2505.16932}, 2025.

\bibitem[Wang et al.(2026)]{wang2026curvature}
S.~Wang, F.~Zhang, J.~Li, D.~Bergemann, and Z.~Yang.
\newblock Why Muon Outperforms Adam: A Curvature Perspective.
\newblock \emph{arXiv:2606.04662}, 2026.

\bibitem[Wen et al.(2026)]{wen2026hyperball}
K.~Wen, X.~Dang, K.~Lyu, T.~Ma, and P.~Liang.
\newblock Fantastic Pretraining Optimizers and Where to Find Them II: Hyperball
  Optimization.
\newblock \emph{arXiv:2606.16899}, 2026.
\end{thebibliography}

\end{document}
